% generated from papers/V-free/paper_V_free.tex -- edit the paper, then rerun assemble.py
\chapter[The free variant]{The free variant and the free critical exponent}\label{ch:V}
\chapterorigin{This chapter is Paper V of the series (\texttt{papers/V-free/}).}
\begin{summary}
Replacing the free abelian monoid $J_S$ of the localized Bost--Connes system by the free
monoid $\F_S^+$ produces a system whose statistics are free-probabilistic and whose
thermodynamics is governed by a new arithmetic invariant. We determine both, and conclude
that the deformation is incompatible with the phase structure it was meant to probe.

The incompatibility begins at the definition: keeping the coefficient algebra $C(Y_S)$ while
freeing the isometries is inconsistent, since orthogonal range projections contradict the
nested ones of $C(Y_S)$, and the universal algebra is zero. The consistent free system is
the Toeplitz--Cuntz algebra of $\F_S^+$ with the norm dynamics; it retains the norms of
$S$ and nothing else --- no symmetry group $G_S$, no obstruction group $\Xi_\beta$. Its
partition function is $Z_{\mathrm{free}}(\beta)=(1-L_S(\beta))^{-1}$ with
$L_S(\beta)=\sum_{\pp\in S}\Ns\pp^{-\beta}$, so the critical point solves $L_S=1$: a
Galton--Watson branching criticality rather than an abscissa of convergence. The range
projections being mutually orthogonal, every $\mathrm{KMS}_\beta$ state forces
$L_S(\beta)\le1$, so \emph{no $\mathrm{KMS}_\beta$ state exists below the threshold}, and
above it there is exactly one --- the same mechanism by which partitioning, as opposed to
nested, range projections restrict the admissible temperatures in the $ax+b$ systems
(Chapter~\ref{ch:II}).

In the Fock--Gibbs state the field operators $\Phi_\pp=\mu_\pp+\mu_\pp^*$ have an exactly
computable law on $[-2,2]$: with $t=\Ns\pp^{-\beta}$ the moments are $m_2=1+t$,
$m_4=2+3t+t^2$ and the free cumulant $\kappa_4=-t(1+t)$ is strictly negative, so
$\Phi_\pp$ is semicircular only in the limit $\beta\to\infty$. The fields are not exactly
free: the obstruction is the hard-core exclusion $e_\pp e_\qq=0$, giving the mixed cumulant
$-t_\pp t_\qq$, and asymptotic freeness holds at exactly that rate. Two natural guesses fail:
$\cN_\pp=\mu_\pp\mu_\pp^*$ is a \emph{projection}, so its law is Bernoulli and not free
Poisson; and Voiculescu's microstates free entropy is undefined for $n\ge2$ variables, the
$\mathrm{KMS}_\beta$ state not being tracial.

The invariant $\bfree(S)$, the root of $L_S=1$, is arithmetically of a different kind from
the classical critical exponent. It is not a tail invariant: over $\Q$, deleting the single
prime $2$ moves it from $1.399433\ldots$, the root of the prime zeta function minus one,
to $1.224388\ldots$; and it obeys the universal bound $\bfree(S)\le\bfree(P_\Q)$. Applied to
the cascade constructions of [Ch.~\ref{ch:I}] the conclusion is sharp and negative: there
$L_S(1)=\infty$ by design, so $\bfree(S)>1$, while the entire transition locus lies in
$(0,1]$. Every classical phase transition occurs at a temperature where the free system has no
equilibrium state at all, and $\bfree$ can be prescribed anywhere in $\bigl(1,\bfree(P_\Q)\bigr)$
without disturbing the cascade. The two exponents carry disjoint information.
\end{summary}

\section{The free system}

\subsection{The naive deformation is inconsistent}

The first idea is to keep the coefficient algebra of the classical system and to free the
isometries: consider the universal $C^*$-algebra $\cA^{\mathrm{naive}}$ generated by a copy
of $C(Y_S)$ and isometries $\mu_\pp$, $\pp\in S$, subject to
\begin{equation}\label{V:eq:naive}
\mu_\pp^*\mu_\qq=\delta_{\pp\qq}1,\qquad \mu_\pp f\mu_\pp^*=\alpha_\pp(f)\quad(f\in C(Y_S)),
\end{equation}
where $\alpha_\pp(f)=(f\circ\pp^{-1})\mathbf 1_{\pp Y_S}$ is the endomorphism of the
classical system [Ch.~\ref{ch:Main}, \S2]. The two halves of \eqref{V:eq:naive} pull in opposite
directions: the first makes the range projections $e_\pp=\mu_\pp\mu_\pp^*$ orthogonal, the
second identifies $e_\pp=\alpha_\pp(1)$ with the image of $\mathbf 1_{\pp Y_S}$, and in
$C(Y_S)$ the range projections are nested, $\mathbf 1_{\pp Y_S}\mathbf 1_{\qq Y_S}=\mathbf 1_{\pp\qq Y_S}\ne0$.

\begin{proposition}[Incompatibility]\label{V:prop:naive}
If $|S|\ge2$ then $\cA^{\mathrm{naive}}=0$.
\end{proposition}

\begin{proof}
Let $\pp\ne\qq$. By \eqref{V:eq:naive}, $e_\pp e_\qq=\mu_\pp(\mu_\pp^*\mu_\qq)\mu_\qq^*=0$,
while $e_\pp e_\qq$ is the image of $\mathbf 1_{\pp Y_S}\mathbf 1_{\qq Y_S}=\mathbf 1_{\pp\qq Y_S}=\alpha_\pp(\mathbf 1_{\qq Y_S})$.
Hence $\mu_\pp\mathbf 1_{\qq Y_S}\mu_\pp^*=0$, and multiplying by $\mu_\pp^*$ on the left and
$\mu_\pp$ on the right, using $\mu_\pp^*\mu_\pp=1$, gives $\mathbf 1_{\qq Y_S}=0$ in
$\cA^{\mathrm{naive}}$. Then $\mu_\qq=e_\qq\mu_\qq=\mathbf 1_{\qq Y_S}\mu_\qq=0$, so
$1=\mu_\qq^*\mu_\qq=0$.
\end{proof}

The classical system has a single isometry per prime with \emph{nested} ranges, and this
is what lets $C(Y_S)$, with its symmetry group $G_S$ and obstruction group $\Xi_\beta$, sit
inside it; a free monoid of isometries has \emph{partitioning} ranges, and cannot coexist
with that coefficient algebra. The consistent free deformation therefore retains, of the
arithmetic of $S$, only the norms.

\subsection{The Toeplitz--Cuntz system}

\begin{definition}\label{V:def:free}
Let $\F_S^+$ be the free monoid on $S$ and $\Ns w=\prod_j\Ns\pp_j$ for
$w=\pp_1\cdots\pp_k$. The \emph{free system} is the Toeplitz--Cuntz algebra
$\cA^{\mathrm{free}}_S=\cT_S$, the universal $C^*$-algebra generated by isometries $\mu_\pp$,
$\pp\in S$, with
\begin{equation}\label{V:eq:rel}
\mu_\pp^*\mu_\qq=\delta_{\pp\qq}1 ,
\end{equation}
with the dynamics $\sigma_t(\mu_w)=\Ns w^{it}\mu_w$. Write $e_\pp=\mu_\pp\mu_\pp^*$,
$\mu_w=\mu_{\pp_1}\cdots\mu_{\pp_k}$, and $L_S(\beta)=\sum_{\pp\in S}\Ns\pp^{-\beta}$.
\end{definition}

The Fock representation on $\ell^2(\F_S^+)$, $\mu_\pp\xi_w=\xi_{\pp w}$, is faithful: for
finite $S$ this is Cuntz's theorem \cite{Cuntz77}, and in general it follows from the
gauge-invariant uniqueness theorem, the Fock representation carrying the gauge action. The diagonal $C^*\{\mu_w\mu_w^*\}$ of $\cT_S$ is commutative and plays the
role of the coefficient algebra; its spectrum is the space of finite and infinite words,
the boundary of the word tree, and \emph{not} $Y_S$.

\begin{proposition}\label{V:prop:Z}
$\sum_{w\in\F_S^+}\Ns w^{-\beta}=\sum_{k\ge0}L_S(\beta)^k=(1-L_S(\beta))^{-1}$, convergent
precisely when $L_S(\beta)<1$.
\end{proposition}

\begin{proof}
Words of length $k$ are indexed by $S^k$ and $\Ns$ is multiplicative, so
$\sum_{|w|=k}\Ns w^{-\beta}=L_S(\beta)^k$.
\end{proof}

\begin{theorem}[The equilibrium states]\label{V:thm:betac}
Let $|S|\ge2$. Then $L_S$ is continuous and strictly decreasing on its domain with
$L_S(0^+)>1$ and $L_S(\infty)=0$, so there is a unique $\bfree(S)$ with $L_S(\bfree)=1$
whenever $L_S$ exceeds $1$ somewhere. Moreover:
\begin{enumerate}
\item every $\mathrm{KMS}_\beta$ state satisfies $L_S(\beta)\le1$; hence for
$\beta<\bfree(S)$ \emph{no $\mathrm{KMS}_\beta$ state exists};
\item for every $\beta$ there is at most one $\mathrm{KMS}_\beta$ state, and it is given
on the dense $*$-algebra spanned by the $\mu_w\mu_v^*$ by
\begin{equation}\label{V:eq:unique}
\varphi(\mu_w\mu_v^*)=\delta_{wv}\,\Ns w^{-\beta};
\end{equation}
\item for $L_S(\beta)<1$ the Fock--Gibbs state
\begin{equation}\label{V:eq:gibbs}
\varphi_\beta(x)=\bigl(1-L_S(\beta)\bigr)\sum_{w\in\F_S^+}\Ns w^{-\beta}\bigl\langle x\,\xi_w,\xi_w\bigr\rangle
\end{equation}
on $\ell^2(\F_S^+)$ is a $\mathrm{KMS}_\beta$ state, hence the unique one;
\item at $\beta=\bfree(S)$ a $\mathrm{KMS}_\beta$ state exists, namely the weak-$*$ limit of
$\varphi_{\beta'}$ as $\beta'\downarrow\bfree$, and it is the unique one.
\end{enumerate}
So the free system has no equilibrium state below $\bfree$ and exactly one at every
$\beta\ge\bfree$.
\end{theorem}

\begin{proof}
(1) By \eqref{V:eq:rel} the $e_\pp$ are mutually orthogonal, so $\sum_{\pp\in F}e_\pp\le1$ for
finite $F$; and the KMS condition gives
$\varphi(e_\pp)=\Ns\pp^{-\beta}\varphi(\mu_\pp^*\mu_\pp)=\Ns\pp^{-\beta}$. Hence
$\sum_{\pp\in F}\Ns\pp^{-\beta}\le1$ for every finite $F$, i.e.\ $L_S(\beta)\le1$, and strict
monotonicity forces $\beta\ge\bfree$.

(2) The elements $\mu_w\mu_v^*$ span a dense $*$-subalgebra, by \eqref{V:eq:rel}. For a
$\mathrm{KMS}_\beta$ state $\varphi$ and words $w,v$, the KMS condition with the entire
element $\mu_w$, $\sigma_{i\beta}(\mu_w)=\Ns w^{-\beta}\mu_w$, gives
$\varphi(\mu_w\mu_v^*)=\Ns w^{-\beta}\varphi(\mu_v^*\mu_w)$. If neither of $v,w$ is a prefix
of the other, $\mu_v^*\mu_w=0$ by \eqref{V:eq:rel}. If $w=vu$ with $u\ne\emptyset$ then
$\mu_v^*\mu_w=\mu_u$, and if $v=wu$ then $\mu_v^*\mu_w=\mu_u^*$; in both cases the value
vanishes, since $\varphi$ is invariant under $\sigma_t$ and $\mu_u^{\pm}$ is an eigenvector
of $\sigma_t$ with eigenvalue $\Ns u^{\pm it}\ne1$. Hence $\varphi(\mu_w\mu_v^*)=0$ for
$w\ne v$, and $\varphi(\mu_w\mu_w^*)=\Ns w^{-\beta}\varphi(\mu_w^*\mu_w)=\Ns w^{-\beta}$.

(3) \eqref{V:eq:gibbs} is a state by Proposition~\ref{V:prop:Z}, and
$\langle\mu_w\mu_v^*\xi_u,\xi_u\rangle\ne0$ forces $u=vu'$ and $wu'=u$, i.e.\ $w=v$, so
$\varphi_\beta(\mu_w\mu_v^*)=\delta_{wv}(1-L_S)\sum_{u'}\Ns(wu')^{-\beta}=\delta_{wv}\Ns w^{-\beta}$,
which is \eqref{V:eq:unique}; the KMS condition for the generators then follows from
\eqref{V:eq:unique} by the computation of (2) read backwards.

(4) The state space is weak-$*$ compact, and a weak-$*$ limit of $\mathrm{KMS}_{\beta'}$
states with $\beta'\to\beta$ is a $\mathrm{KMS}_\beta$ state \cite[Prop.~5.3.25]{BR};
the limit exists because \eqref{V:eq:unique} converges as $\beta'\to\bfree$.
\end{proof}

\begin{remark}[Branching versus convergence]\label{V:rem:branch}
Classically the partition function is the Euler product
$\prod_\pp(1-\Ns\pp^{-\beta})^{-1}$ and the critical exponent is the abscissa of convergence
of $L_S$. Freely, $\F_S^+$ has exponentially many words, $L_S(\beta)$ is the branching number
of the weighted word tree, and the transition is the Galton--Watson criticality $L_S=1$: a
tree of branching number $>1$ has too many words to normalize. The same orthogonality that
yields $L_S(\beta)\le1$ produces the constraint $\beta=1$ on the affine ($ax+b$) system,
whose Cuntz relations are finite partitions (Chapter~\ref{ch:II}). Nested range projections leave
every temperature admissible; partitioning ones do not. And the classical phase structure
does not survive at all: by Proposition~\ref{V:prop:naive} the free system has no coefficient
algebra $C(Y_S)$, hence no symmetry group $G_S$, no obstruction group $\Xi_\beta$, and, by
Theorem~\ref{V:thm:betac}(2), a unique equilibrium state wherever it has one. Everything
below is computed in the Fock--Gibbs state.
\end{remark}

\section{The law of the fields}

Fix $\pp$, write $t=\Ns\pp^{-\beta}\in(0,1)$, $\Phi=\mu_\pp+\mu_\pp^*$, $e=e_\pp$.

\begin{theorem}[Moments and cumulants]\label{V:thm:moments}
In $\varphi_\beta$,
\[
m_1=m_3=0,\qquad m_2=1+t,\qquad m_4=2+3t+t^2,\qquad \kappa_4=m_4-2m_2^2=-t(1+t)<0 ,
\]
and the law of $\Phi_\pp$ is
\[
d\rho_t(x)=(1-t)\sum_{n\ge0}t^n\,U_n(x/2)^2\,\frac{\sqrt{4-x^2}}{2\pi}\,dx ,
\]
supported on $[-2,2]$ for every $\beta$, with $U_n$ the Chebyshev polynomials of the second
kind. In particular $\Phi_\pp$ is not semicircular for $t>0$.
\end{theorem}

\begin{proof}
Only gauge-invariant words survive, so odd moments vanish and
$m_2=\varphi(\mu\mu^*+\mu^*\mu)=t+1$. For $m_4$, using $\mu^*\mu=1$ and
$\varphi(\mu x\mu^*)=t\varphi(x)$, the six balanced words give
$\varphi(\mu\mu\mu^*\mu^*)=t^2$, $\varphi(\mu\mu^*\mu\mu^*)=\varphi(e)=t$,
$\varphi(\mu\mu^*\mu^*\mu)=t$, $\varphi(\mu^*\mu\mu\mu^*)=t$,
$\varphi(\mu^*\mu\mu^*\mu)=1$, $\varphi(\mu^*\mu^*\mu\mu)=1$, summing to $t^2+3t+2$. For the
density, note that $C^*(\mu_\pp)$ is the Toeplitz algebra, and that on the Fock space
$\ell^2(\F_S^+)=\bigoplus_{u}\ell^2(\pp^{\N}u)$, the sum over words $u$ not beginning with
$\pp$, the operator $\Phi$ is the Jacobi matrix with unit off-diagonal and zero diagonal on
each summand $\ell^2(\N)$, whose spectral measure at the $n$-th basis vector is
$U_n(x/2)^2\sqrt{4-x^2}\,dx/2\pi$; grouping the terms of \eqref{V:eq:gibbs} by $u$, the law of
$\Phi$ is the geometric mixture with weights $(1-t)t^n$, since the factor
$\sum_u\Ns u^{-\beta}$ over words not beginning with $\pp$ is $(1-t)/(1-L_S(\beta))$. The
relation $\kappa_4=m_4-2m_2^2$ for a centred variable is the moment--cumulant formula
\cite{NS}.
\end{proof}

\begin{corollary}[Semicircular limit]\label{V:cor:semicircle}
As $\beta\to\infty$, $\rho_t$ converges weakly to the Wigner semicircular law of radius $2$
\cite{VDN};
equivalently $\kappa_2\to1$, all free cumulants of order $\ge3$ vanish, and the even moments
converge to the Catalan numbers $1,2,5,14,42,\dots$
\end{corollary}

\begin{theorem}[The number operators are projections]\label{V:thm:bernoulli}
$\cN_\pp=e_\pp$ satisfies $e_\pp^2=e_\pp$, so
$\mathrm{law}(\cN_\pp)=(1-t)\delta_0+t\delta_1$ in $\varphi_\beta$. This is a Bernoulli law,
not a free Poisson law: a free Poisson variable is a projection only in the degenerate cases
$\delta_0$ and $\delta_1$.
\end{theorem}

\begin{proof}
$\mu_\pp$ is an isometry, so $\varphi(e_\pp^n)=\varphi(e_\pp)=t$ for every $n\ge1$, which
determines the stated law; its spectrum is $\{0,1\}$, whereas a nondegenerate
Marchenko--Pastur law has an absolutely continuous part.
\end{proof}

\section{Asymptotic, not exact, freeness}

\begin{theorem}[The obstruction]\label{V:thm:notfree}
For $\pp\ne\qq$ the fields $\Phi_\pp,\Phi_\qq$ are not free in $\varphi_\beta$: while
$\varphi(\Phi_\pp\Phi_\qq)=0$ and $\varphi(\Phi_\pp\Phi_\qq\Phi_\pp\Phi_\qq)=0$, one has
\begin{equation}\label{V:eq:mixed}
\varphi\bigl(\Phi_\pp^2\Phi_\qq^2\bigr)-\varphi\bigl(\Phi_\pp^2\bigr)\varphi\bigl(\Phi_\qq^2\bigr)=-\,t_\pp t_\qq .
\end{equation}
\end{theorem}

\begin{proof}
Gauge invariance requires each letter to appear equally often with and without adjoint, so
$\varphi(\Phi_\pp\Phi_\qq)=0$. For $\Phi_\pp^2\Phi_\qq^2$ the surviving terms give
$\varphi\bigl((e_\pp+1)(e_\qq+1)\bigr)=\varphi(e_\pp e_\qq)+t_\pp+t_\qq+1$, and
$e_\pp e_\qq=\mu_\pp(\mu_\pp^*\mu_\qq)\mu_\qq^*=0$ by \eqref{V:eq:rel}; subtract
$(1+t_\pp)(1+t_\qq)$. For the alternating moment the four balanced terms are
$\mu_{\pp\qq}\mu_{\qq\pp}^*$, $\mu_\pp\mu_{\pp\qq}^*\mu_\qq$,
$\mu_\pp^*\mu_\qq\mu_\pp\mu_\qq^*$ and $\mu_{\qq\pp}^*\mu_{\pp\qq}$; the second and third
vanish by $\mu_\pp^*\mu_\qq=0$, the fourth is $\mu_{\qq\pp}^*\mu_{\pp\qq}=0$, and the
first is $\varphi(\mu_w\mu_v^*)$ with $w\ne v$, hence $0$ by \eqref{V:eq:unique}.
\end{proof}

\begin{corollary}[Asymptotic freeness]\label{V:cor:asympfree}
Let $(\pp_k)\subset S$ with $\Ns\pp_k\to\infty$ and $\beta>\bfree(S)$. Every mixed free
cumulant of $\{\Phi_{\pp_k}\}$ is $O\bigl(\prod_jt_{\pp_{i_j}}\bigr)$, hence
$O\bigl(\min_j\Ns\pp_{i_j}^{-\beta}\bigr)$, and tends to $0$; the family is asymptotically
free and by \eqref{V:eq:mixed} the rate $t_\pp t_\qq$ is attained.
\end{corollary}

\begin{proof}
By \eqref{V:eq:unique} a mixed moment $\varphi(\Phi_{\pp_{i_1}}\cdots\Phi_{\pp_{i_n}})$ is a sum,
over the choices of $\mu$ or $\mu^*$ in each factor, of terms $\varphi(\mu_w\mu_v^*)$ after
reduction by \eqref{V:eq:rel}; a term survives only if it reduces to $\mu_w\mu_w^*$, and then
equals $\Ns w^{-\beta}=\prod t$ over the letters of $w$ with multiplicity. Every letter
distinct from those reducing to the identity contributes a factor $t_\pp$, so the moment
differs from the corresponding moment of free variables with the same one-variable laws by
terms each carrying at least one factor $t_\pp$ per additional distinct letter. Möbius
inversion on $\NC(n)$ \cite{NS}, a finite linear combination of products of moments,
preserves this order of magnitude, and each $t<1$.
\end{proof}

\begin{remark}[Hard-core exclusion]\label{V:rem:hardcore}
Free variables would satisfy $\varphi(e_\pp e_\qq)=t_\pp t_\qq$; the true value is $0$,
because no vector lies simultaneously in the ranges of $\mu_\pp$ and $\mu_\qq$. The
discrepancy is exactly $-t_\pp t_\qq$, and as the norms grow the ranges become
$\varphi$-small and the exclusion invisible. This is the free counterpart of the classical
asymptotic independence of distinct primes for the geometric measures $\pi_\beta$.
\end{remark}

\begin{theorem}[Free entropy is unavailable]\label{V:thm:nochi}
$\varphi_\beta$ is not tracial: $\varphi_\beta(\mu_\pp\mu_\pp^*)=t_\pp<1=\varphi_\beta(\mu_\pp^*\mu_\pp)$.
Hence Voiculescu's microstates free entropy $\chi(\Phi_{\pp_1},\dots,\Phi_{\pp_n})$
\cite{Voi94} and the free Fisher information are undefined for $n\ge2$, both requiring a
tracial $W^*$-probability space. For $n=1$ the entropy is defined, the distribution of a single
self-adjoint element being a measure, and equals
$\chi(\Phi_\pp)=\iint\log|x-y|\,d\rho_t(x)d\rho_t(y)+\tfrac34+\tfrac12\log2\pi$; it is finite
and continuous in $t$, and since $t_\pp<L_S(\bfree)=1$ strictly, it remains finite as
$\beta\downarrow\bfree$. There is no logarithmic singularity of the free entropy at the
critical point.
\end{theorem}

\begin{proof}
The displayed inequality is the KMS condition. Both invariants are defined for $n$-tuples in
a tracial von Neumann algebra, and no such structure is available. For $n=1$, $\rho_t$ is
supported in $[-2,2]$ with bounded density for $t$ bounded away from $1$, so the logarithmic
energy is finite and continuous.
\end{proof}

\begin{remark}[Two things called free]\label{V:rem:freeenergy}
The quantity $-\log(1-L_S(\beta))=\log Z_{\mathrm{free}}(\beta)$ does diverge as
$\beta\downarrow\bfree$, by Proposition~\ref{V:prop:Z}. That is the free \emph{energy}, a
thermodynamic quantity attached to the whole algebra, not Voiculescu's free \emph{entropy},
an invariant of an $n$-tuple of noncommutative random variables. The coincidence of names is
unfortunate; an asymptotic of the form $\chi_\beta\sim\frac n2\log(1-L_S(\beta))$ conflates
them, its right side being the free energy and its left side undefined for $n\ge2$ and
bounded for $n=1$.
\end{remark}

\section{The free critical exponent}

\begin{definition}\label{V:def:bfree}
Let $\bcl(S)=\inf\{\beta:L_S(\beta)<\infty\}$ be the abscissa of convergence, the classical
critical exponent of [Ch.~\ref{ch:Main}], and
$\lambda_S=\lim_{\beta\downarrow\bcl(S)}L_S(\beta)\in(0,\infty]$. When $\lambda_S>1$ let
$\bfree(S)$ be the root of $L_S=1$; when $\lambda_S\le1$ set $\bfree(S)=\bcl(S)$.
\end{definition}

\begin{theorem}[Trichotomy]\label{V:thm:trich}
For infinite $S$ exactly one of the following holds.
\begin{enumerate}
\item[\textup{(S)}] $\lambda_S>1$: there is a unique $\bfree(S)\in(\bcl(S),\infty)$, the
admissible set is $[\bfree,\infty)$, and $Z_{\mathrm{free}}\to\infty$ as
$\beta\downarrow\bfree$.
\item[\textup{(C)}] $\lambda_S=1$: $\bfree=\bcl$, the admissible set is $(\bcl,\infty)$, and
$Z_{\mathrm{free}}\to\infty$ at $\bcl$.
\item[\textup{(B)}] $\lambda_S<1$: $\bfree=\bcl$, the admissible set is $(\bcl,\infty)$, and
$Z_{\mathrm{free}}$ is \emph{bounded} there, with supremum $(1-\lambda_S)^{-1}<\infty$: the
free phase is closed.
\end{enumerate}
For every $\alpha\in(0,1)$ and every $\lambda\in(0,\infty]$ there is $S\subseteq P_\Q$ with
$\bcl(S)=\alpha$ and $\lambda_S=\lambda$, so all three occur.
\end{theorem}

\begin{proof}
Monotonicity and continuity of a Dirichlet series on its half-plane of convergence give the
trichotomy, using Theorem~\ref{V:thm:betac}(1) for the admissible sets. For the realization,
fix $\alpha$. A set $T$ with $|T\cap[1,x]|\asymp x^\alpha$, obtained greedily as in
[Ch.~\ref{ch:I}, Lem.~5.4], has $\bcl(T)=\alpha$ and, by Abel summation, $L_T(\alpha)=\infty$,
giving $\lambda=\infty$. For $\lambda<\infty$ take instead
$|T\cap[1,x]|\asymp x^\alpha/\log^2x$; then $\bcl(T)=\alpha$ still while the $m$-th element
satisfies $\ell_m\asymp(m\log^2m)^{1/\alpha}$ and $\ell_m^{-\alpha}\asymp(m\log^2m)^{-1}$, so
$L_T(\alpha)<\infty$. Discarding an initial segment of $T$ lowers $L_T(\alpha)$ below
$\lambda$ without changing $\bcl$; then adjoin primes $p_1<p_2<\cdots$ outside $T$ chosen
greedily, $p_k$ so large that $p_k^{-\alpha}$ is at most half of the remaining deficit
$\lambda-L(\alpha)$ of the set built so far. The deficit tends to $0$, the adjoined set has
counting function $O(\log x)$, hence does not change $\bcl$, and the resulting $S$ has
$L_S(\alpha)=\lambda$ exactly.
\end{proof}

\begin{proposition}[Not a tail invariant]\label{V:prop:tail}
$\bcl(S)$ is unchanged by adding or deleting finitely many primes and depends only on the
counting function, through $\bcl(S)=\limsup_m\log m/\log\Ns\pp_m$. By contrast $\bfree$ is
monotone in $S$ and strictly changed by deleting a single prime in regime (S); two sets with
the same counting asymptotics can have different $\bfree$. Moreover, over $\Q$,
\[
\bfree(S)\le\bfree(P_\Q)=1.399433335\ldots ,
\]
the root of the prime zeta function minus one, with equality only for $S=P_\Q$.
\end{proposition}

\begin{proof}
The first is the standard formula for the convergence exponent. The rest follows from strict
monotonicity of $L_S$ in $S$ and in $\beta$.
\end{proof}

\begin{remark}[Computed values]\label{V:rem:values}
Computed from the primes below $x=2\cdot10^7$ with an analytic tail correction (code in
\texttt{code/verify\_series.py}): substituting $t=e^u$ gives
$\sum_{\pp>x}\Ns\pp^{-\beta}\approx\int_{\log x}^\infty e^{(1-\beta)u}u^{-k}du$, the
exponential integral $E_1((\beta-1)\log x)$ for $k=1$ and, for the Hardy--Littlewood
densities $2C_2/\log^2t$ with $k=2$, a quantity finite at $\beta=1$ with value
$2C_2/\log x=0.0785$:
\[
\begin{array}{lcc}
S & L_S(1^+) & \bfree(S)\\\hline
\text{all primes} & \infty & 1.399433\\
\text{all primes except }2 & \infty & 1.224388\\
\text{all primes except }2,3 & \infty & 1.154890\\
\{p>10^6\} & \infty & 1.019162\\
\text{Sophie Germain} & 1.535 & 1.2114\\
\text{twin primes (lower)} & 1.059 & 1.0089
\end{array}
\]
The first four rows are correct to the digits shown; the last two depend on the tail
model in their third decimal. Removing the single prime $2$ moves $\bfree$ by $0.175$,
quantifying Proposition~\ref{V:prop:tail}; nothing comparable can happen to $\bcl$, which is
$1$ for the first four rows. The twin row should be read with care: the bare sum over the
computed range is $0.9805$ and the Hardy--Littlewood tail contributes $0.0785$, so the set
is supercritical but only just, and the conclusion rests on the tail estimate rather than
on a computation (Brun's constant $B_2=1.9022\ldots$ gives the same picture:
$\sum_{p\ \mathrm{twin}}p^{-1}=\tfrac12\bigl(B_2+\sum2/(p(p+2))\bigr)\approx1.06$).
Whether $\sum_{p\ \mathrm{twin}}p^{-1}$ exceeds $1$ is an arithmetic question we do not
settle; if it does not, the twin set is an instance of regime (B).
\end{remark}

\section{The free deformation annihilates the cascade}

\begin{theorem}[Annihilation]\label{V:thm:annih}
Let $S$ be any of the cascade or Cantor sets of [Ch.~\ref{ch:I}, Thm.~5.1, Cor.~6.2]. Then
$\bcl(S)=1$ and $\lambda_S=\infty$, so $\bfree(S)>1$ strictly; the transition locus of the
classical system lies in $(0,1]$; and by Theorem~\ref{V:thm:betac}(1) the free system has no
$\mathrm{KMS}_\beta$ state for $\beta<\bfree(S)$. Consequently every classical phase
transition occurs at a temperature at which the free system has no equilibrium state; on the
whole admissible range of the free system the classical one has $\zeta_{K,S}(\beta)<\infty$
and $\Xi_\beta=\widehat{G_S}$, i.e.\ is in its Gibbs regime with all its phase transitions
behind it; and the free system has a single transition, at $\bfree(S)$, from no
equilibrium state to exactly one.
\end{theorem}

\begin{proof}
The abundant base of [Ch.~\ref{ch:I}, Lem.~5.4(1)] has $\sum_{\ell\in P_0}\ell^{-1}=\infty$, so
$L_S(1)=\infty$, $\bcl(S)=1$ and $\lambda_S=\infty$; Theorem~\ref{V:thm:trich}(S) gives
$\bfree(S)>1$. By [Ch.~\ref{ch:I}, Thm.~5.1(3)] the thresholds lie in $(0,1)$ and
$\Xi_\beta=\widehat{G_S}$ for $\beta>1$.
\end{proof}

\begin{theorem}[The two exponents are independent]\label{V:thm:tune}
Let $S$ be a cascade set and $\gamma\in\bigl(1,\bfree(P_\Q)\bigr)$. There is $S'$ differing
from $S$ by finitely many primes --- hence with the same chain $\{\Xi_\beta\}$ and the same
transition locus --- with $\bfree(S')$ within any prescribed $\varepsilon$ of $\gamma$. The
interval is optimal: $\bfree(S')>1$ always by Theorem~\ref{V:thm:annih} and
$\bfree(S')<\bfree(P_\Q)$ by Proposition~\ref{V:prop:tail}.
\end{theorem}

\begin{proof}
Adding or deleting finitely many primes leaves $\Xi_\beta$ unchanged, a finite set
contributing a finite amount to every series in \eqref{V:eq:rel} and the computation of
$\Xi_\beta$ in [Ch.~\ref{ch:I}, Thm.~5.1] being stated modulo finite exceptional sets.
Deleting all primes of $S$ below $z$ gives
$L_{S'}(\beta)\le\sum_{\ell>z}\ell^{-\beta}\to0$ for fixed $\beta>1$ as $z\to\infty$, so
$\bfree(S')$ can be pushed arbitrarily close to $1$; adjoining all primes below $N$ gives
$L_{S'}\ge\sum_{p\le N}p^{-\beta}$, whose root increases to $\bfree(P_\Q)$. Adding one prime
at a time moves the root by a controlled amount, giving every intermediate value up to
$\varepsilon$.
\end{proof}

\begin{remark}[Numerical illustration]\label{V:rem:tuning}
Restricting to primes above $z$ gives the upper bounds $1.056773$ at $z=10^2$, $1.028710$ at
$z=10^4$, $1.019162$ at $z=10^6$ and $1.014372$ at $z=10^8$ (the last from the tail
integral alone), tending to $1$; adjoining all
primes below $N$ gives roots $1.147338$ at $N=10$, $1.356444$ at $N=10^2$, $1.388390$ at
$N=10^3$, $1.396110$ at $N=10^4$, $1.399069$ at $N=10^6$ and $1.399307$ at $N=10^7$,
increasing to $\bfree(P_\Q)=1.399433$ but never attaining it.
\end{remark}

\begin{remark}[What the deformation is good for]\label{V:rem:moral}
The free variant was proposed as a route to importing free probability into the arithmetic
phase structure. Theorem~\ref{V:thm:annih} shows the two are incompatible in a precise sense:
the phenomena classified by $\Xi_\beta$ live at high temperature, where the partition
function diverges, and it is exactly that region which the free monoid renders inaccessible,
because exponential word growth prevents normalization once the branching number exceeds
$1$. The classical system is interesting at high temperature and trivial at low; the free
system exists only at low temperature and is trivial there. Theorem~\ref{V:thm:tune} confirms
that the two exponents carry disjoint information rather than being two views of one
structure.

What the deformation does produce is $\bfree$ itself, a well-defined arithmetic invariant of
a set of primes: computable, universally bounded over $\Q$, and --- unlike every exponent in
this series --- sensitive to individual norms rather than only to a counting function.
\end{remark}

\section{Assessment}

\[
\begin{array}{lll}
\text{Statement} & \text{Status} & \text{Reference}\\\hline
C(Y_S)\ \text{survives the deformation} & \text{false: the universal algebra is }0 & \text{Prop.~\ref{V:prop:naive}}\\
Z_{\mathrm{free}}=(1-L_S)^{-1},\ L_S(\bfree)=1 & \text{true} & \text{Prop.~\ref{V:prop:Z}, Thm.~\ref{V:thm:betac}}\\
\text{two regimes across }\bfree & \text{false: empty below} & \text{Thm.~\ref{V:thm:betac}(1)}\\
\text{KMS}_\beta\text{ state unique for }\beta\ge\bfree & \text{true} & \text{Thm.~\ref{V:thm:betac}(2)}\\
\Phi_\pp\ \text{semicircular} & \text{only as }\beta\to\infty & \text{Thm.~2.1, Cor.~2.2}\\
\cN_\pp\ \text{free Poisson} & \text{false: a projection} & \text{Thm.~2.3}\\
\{\Phi_\pp\}\ \text{exactly free} & \text{false} & \text{Thm.~3.1}\\
\{\Phi_\pp\}\ \text{asymptotically free} & \text{true} & \text{Cor.~3.2}\\
\chi(\Phi_{\pp_1},\dots,\Phi_{\pp_n}),\ n\ge2 & \text{undefined} & \text{Thm.~3.4}\\
\bfree\ \text{a tail invariant} & \text{false} & \text{Prop.~4.3}\\
\text{cascade survives the deformation} & \text{false} & \text{Thm.~5.1}
\end{array}
\]

Three questions remain. First, whether some intermediate deformation --- a coefficient
algebra adapted to the word tree, carrying an action of a quotient of $G_S$ --- can retain
part of the classical phase structure while keeping partitioning ranges; Proposition~\ref{V:prop:naive}
rules out the obvious candidate, not every candidate. Second, $\varphi_\beta$ admits no
trace, but the associated von Neumann algebra has a canonical trace on its continuous core;
whether a free entropy computed there carries arithmetic information is open, and by the
pattern of Chapter~\ref{ch:III} we expect not. Third, $\bfree(S)$ is defined
for any set of primes and we have not found an interpretation of it outside this system ---
for instance as a threshold for a branching process on the multiplicative semigroup generated
by $S$ --- nor decided the regime of the twin primes, which by Remark~\ref{V:rem:values}
reduces to comparing a Brun-type constant with $1$.

